% ===== tr-08 · b14 · tr08b：第 5 章 p28--30（原文页 28--30）=====
% [原文 p.28]
这样的 $\phi$ 称为「势函数」（potential function）。

给定 $X_K$ 的模型 $\mathcal X$，$L$ 的模型 $\mathcal L$ 是线丛 $\mathcal L\to\mathcal X$，且 $\mathcal L|_X=L$。对这组数据我们可以联系 $L$ 上唯一的度量 $\|\cdot\|_{\mathcal L}$，它具有如下性质：若 $s$ 是 $\mathcal L$ 在开集 $\mathcal U\subset\mathcal X$ 上的处处非零局部截面，则 $\|s\|_{\mathcal L}\equiv 1$ 在 $\mathcal U\cap X_K$ 上成立。这是良定义的，因为任意两个这样的截面相差一个可逆函数的乘法，而该函数的 NA 绝对值等于常数 1。这一构造可以推广到 $\Q$-线丛，如此得到的度量称为模型度量。它们在连续度量中是稠密的。

为了看出复几何解释，我们设想代数曲线上某个 snc 模型 $\mathcal X$ 上的线丛 $\mathcal L$。给 $\mathcal L$ 配备任意光滑 Hermite 度量 $h$。给定 $\mathcal L$ 的局部截面 $s$，与式 (5.13) 相容的做法是考虑 $X_t$ 上的局部函数
\[
z\mapsto |s(z)|^{1/|\log|t||}_h.
\]
取 $t\to 0$ 的极限，我们恰好得到模型度量。注意 Hermite 度量选取的歧义在极限中被消除。

% [Tutor Note: 本页起原文用花体 $\mathcal L$（模型线丛，活在模型 $\mathcal X$ 上）与正体 $L$（$X_K$ 上的线丛）严格区分；文字层统一抽成 "L"，若照抄会丢掉这层区分。$\|s\|_{\mathcal L}\equiv 1$ 的 $\equiv$ 表示「恒等于常数 1」，不要写成 $=1$。]

Kähler 几何中一个首要的概念是度量的正性，通常用势的多重次调和（psh）性质来表述。上述讨论提示，在 NA 情形（即 $t\to 0$）下，这一概念应当可以表达为线丛的某种数值性质。

\begin{definition}{5.5.1}[Semipositivity I]
[4, Thm. 2.17] 设 $\|\cdot\|$ 是 $L$ 上的模型度量，联系于 $X_K$ 的模型 $\mathcal X$ 上的 $\Q$-线丛 $\mathcal L$。则
\begin{itemize}
\item 度量 $\|\|$ 是半正模型度量，当且仅当 $\mathcal L$ 是 nef 的，即对任何包含在 $\mathcal X_0$ 中的射影曲线 $C$ 都有 $\mathcal L\cdot C\ge 0$；
\item 连续度量 $\|\|e^{-\phi}$ 是半正的，当且仅当它是 $X_K^{an}$ 上一列半正模型度量的一致极限。
\end{itemize}
\end{definition}

\begin{remark}{5.10}
用「nef」代替「ample」的好处在于：若我们把模型进一步爆破，模型线丛的拉回仍保持 nef，但丰沛性会丢失。
\end{remark}

% [Tutor Note: 原文 Definition 5.1 的引用写作 [4, Thm. 2.17]，Remark 10 紧接其后；本章「注」按原文编号（Remark 9、10、11…）套用 $\{5.n\}$，故这里是 $\begin{remark}{5.10}$。nef = numerically effective（数值有效），丰沛 = ample，二者在模型爆破下的行为不同，这正是原文用 nef 的理由。]

在 Kähler 几何中，psh 函数的定义在复坐标卡中是局部的。由于对偶相交复形是单纯对象，人们会期待 NA 的类比概念与凸函数有关。这一直觉部分成立：

\begin{proposition}{5.5.2}[6, Prop 5.9]
设 $\mathcal X$ 是 $X_K$ 的 snc 模型，$\mathcal L\to\mathcal X$ 是 $L\to X$ 的模型线丛，其相联系的闭 $(1,1)$-形式为 $\theta$。则任何连续 $\theta$-psh 函数在 $\Delta_{\mathcal X}\subset X_K^{an}$ 的任何面上的限制都是凸的。
\end{proposition}

% [Tutor Note: 命题 5.2 的两处引用位置原文写作 "Proposition 5.2. [6, Prop 5.9]"，即文献编号紧跟在命题号之后，而不是像 Definition 5.1 那样放在括号里；译文按原文语序保留。$\theta$-psh 指「相对于 $\theta$ 的多重次调和」，不要拆成两个概念。]

图景是：一般的 $\theta$-psh 函数在 $\Delta_{\mathcal X}$ 的各面上定义凸函数，而其中 $\theta$-psh 模型函数给出网格越来越细的分段仿射逼近。

% [原文 p.29]
\subsection{用 Fubini--Study 度量逼近}
Kähler 几何中的一个基本结果是：整类（integral class）中的任何 Kähler 度量都可以由与射影嵌入相联系的 Fubini--Study 度量逼近。通常的 Fubini--Study 度量依赖于整体截面的 $\C$-向量空间上一个 Hermite 内积的选取，而 NA 类比则依赖于 $K$-向量空间 $V=H^0(X_K,mL)$（$m\gg 1$）上的一个 NA 范数，并满足超度量性质 $\|x+y\|_V\le\max\{\|x\|_V,\|y\|_V\}$。在我们的情形 $K=\C((t))$，可以选取 $V$ 的 $K$-基 $s_0,s_1,\dots s_N$，使得
\[
\|a_0s_0+\dots+a_Ns_N\|_V=\max\{|a_0|\|s_0\|_V,\dots,|a_N|\|s_N\|_V\},\qquad \forall a_i\in K.
\]
线丛 $L\to X_K$ 上的 NA Fubini--Study 度量可以定义为
\[
\|s\|_{FS}(x)=\inf_{\tilde s\in V,\ \tilde s(x)=s^{\otimes m}(x)}\|\tilde s\|_V^{1/m},\qquad \forall x\in X_K^{an}.
\]
具体地说，在正交基下并在局部平凡化中写出，有
\begin{equation}
\|s\|_{FS}(x)=\frac{|s(x)|}{\max_j\{|s_j(x)|/\|s_j\|_V\}^{1/m}},\qquad \forall x\in X^{an}. \tag{5.14}
\end{equation}

% [Tutor Note: (14) 的分母里，幂 $1/m$ 作用在 $\max_j\{\cdot\}$ 整体上（即 $(\max_j |s_j(x)|/\|s_j\|_V)^{1/m}$），不是只作用在 $\|s_j\|_V$ 上；文字层抽成 "1/m" 紧跟花括号，务必按页图确认结合方式。另外 $\inf$ 的下标是 $\tilde s\in V$ 且 $\tilde s(x)=s^{\otimes m}(x)$ 两个条件，其中 $s^{\otimes m}$ 是 $L$ 的 $m$ 次张量幂。]

Fubini--Study 逼近定理的一个 NA 类比给出了半正度量的另一种视角：

\begin{proposition}{5.5.3}[Semipositivity II]
[18] 设 $L\to X_K$ 是丰沛的。则 $L$ 上的连续度量是半正的，当且仅当它可以写成一族 Fubini--Study 度量的一致极限。
\end{proposition}

\begin{remark}{5.11}
在逼近定理中，我们可以假设 $s_i$ 是有限 Laurent 多项式。
\end{remark}

% [Tutor Note: 「整类（integral class）」指 $H^2$ 中落在 $H^2(X,\Z)$ 像里的 K\"ahler 类，这里译为「整类」并附原文；不要与「整性」或「整模型」混淆。Remark 11 的 "finite Laurent polynomials" 指 $s_i$ 是 $t$ 的有限 Laurent 多项式，这是后文 $\|s\|_{FS,t}$ 能对有限 $t$ 求值的前提。]

为得到复几何解释，我们照例假设 $X_K$ 是代数退化族 $X$ 的基变换，并带丰沛的极化线丛 $L$。对任意给定的 NA Fubini--Study 度量式 (5.14)，我们可以联系 $(X_t,L)$ 上一族 Fubini--Study 度量：
\begin{equation}
\|s\|_{FS,t}(z)=\frac{|s(z)|}{\Big\{\Big\{\sum_j|s_j(z,t)|^2|t|^{2\log\|s_j\|_V}\Big\}^{1/2m}},\qquad \forall z\in X_t. \tag{5.15}
\end{equation}
这里 $s_j$ 对有限 $t$ 有意义，因为它们被取为 $t$ 的有限 Laurent 多项式。当 $t\to 0$ 时，这些 Fubini--Study 度量收敛到 NA 类比；更确切地说，$\|s\|_{FS,t}^{1/|\log|t||}$ 在 $X\sqcup X_K^{an}$ 上的混合拓扑下收敛到式 (5.14)。

% [Tutor Note: 原文 p29 的 (15) 分母排印为「{{∑_j…}」，即两个左花括号只对应一个右花括号，外层左括号未闭合（已用 800 dpi 裁图确认：p29，x=244--380、y=538--572）。译文按原样视觉复现（\Big\{\Big\{ 与单个 \Big\}^，定界符不要求 \left/\right 配对，不影响编译）；通常应理解为 $\big\{\sum_j|s_j(z,t)|^2|t|^{2\log\|s_j\|_V}\big\}^{1/2m}$，即分母是求和式的 $1/2m$ 次幂。另注意 $|t|$ 的指数为 $2\log\|s_j\|_V$（两个因子相乘），不是 $2\lVert s_j\rVert_V^2$。]

% [原文 p.30]
\subsection{NA Monge--Ampère 测度}
NA Monge--Ampère 测度 [11] 是以某种反直觉的方式通过相交理论定义的。作为动机，我们考虑复解析情形：代数曲线上的 snc 模型 $\mathcal X$，配备 Hermite 线丛 $(\mathcal L,h)$，其曲率形式 $\theta$ 属于类 $c_1(\mathcal L)$。于是 $\theta^n$ 定义了 $X_t$ 上一族 $n$-形式，使得 $\int_{X_t}\theta^n$ 等于相交数 $(L^n)$。当我们把 $X_t$ 视为收敛到对偶相交复形 $\Delta_{\mathcal X}$（参看第 5.3.1 节）时，问题就是描述这些 $n$-形式的极限。记 $\mathcal X_0=\sum_{i\in I}b_iE_i$。回忆：$X_t$ 上对应于对偶相交复形中各面的区域，从代数角度看只是 $E_J$ 的小邻域。因此 $\theta^n|_{X_t}$ 的极限只能支撑在 $\Delta_{\mathcal X}$ 的顶点上，而顶点对应于分支 $E_i$。集中在顶点处的 delta 质量分别为
\[
b_i\int_{E_i}\theta^n=b_i\mathcal L^n\cdot E_i,
\]
其中 $b_i$ 的出现源于叶层（sheets）的重数。令人放心的是，
\[
\sum_i b_i\mathcal L^n\cdot E_i=(L^n)
\]
给出正确的总质量。

回到 NA 情形。给定 $L\to X_K$ 的模型 $\Q$-线丛 $\mathcal L\to\mathcal X$，记 $\mathcal X_0=\sum_i b_iE_i$，并把与 $E_i$ 相联系的除子型点记作 $q_i$。于是我们可以把模型度量 $\|\cdot\|_{\mathcal L}$ 的 NA Monge--Ampère 测度定义为如下支撑在 $q_i\in X_K^{an}$ 上的带符号原子测度：
\[
MA(\|\cdot\|_{\mathcal L})=\sum_{E_i}b_i(\mathcal L^n\cdot E_i)\delta_{q_i}.
\]
这一定义与线丛的拉回通过投影公式相容，并保证总质量是相交数 $(L^n)$。若 $\|\cdot\|_{\mathcal L}$ 还是半正的，则这些相交数非负，于是 $MA(\|\cdot\|_{\mathcal L})$ 是一个测度。

% [Tutor Note: $MA(\|\cdot\|_{\mathcal L})$ 中的 "MA" 是数学斜体（$MA$），不是正体缩写；求和指标原文写的是 $E_i$（不是 $i$），而系数是 $b_i(\mathcal L^n\cdot E_i)$，两者指标必须一致。$\delta_{q_i}$ 是支撑在除子型点 $q_i$ 处的 Dirac 质量；$b_i$ 是 $\mathcal X_0$ 中分支 $E_i$ 的重数，正是它使总量回到 $(L^n)$。]

NA MA 测度的理论与复 MA 测度 [4][5] 高度相似：
\begin{itemize}
\item 在复解析世界，人们首先对光滑势定义复 MA；Kähler 类中一般的连续半正势是光滑势的一致极限，其复 MA 测度则由 $C^0$-收敛下的弱连续性确定；
\item 在 NA 世界，人们首先对模型度量定义 NA MA 测度；$L$ 上一般的连续半正度量是一列连续半正模型度量的一致极限 [6, Cor. 8.8]，其 NA MA 测度可以定义为该序列的 NA MA 测度的唯一极限 Radon 测度 [4, Cor. 3.5]。
\end{itemize}

它们的主要差别在于 NA MA 测度具有高度非局部的面貌。Vilsmeier [76] 的近期结果提供了更具体的视角：

% [Tutor Note: 两个 bullet 的对照结构是「先定义光滑/模型度量的 MA，再用一致极限 + 弱连续性/唯一极限 Radon 测度延拓」；中文不要把它压缩成一句，否则会丢掉「先定模型、再取极限」这一 NA 理论的关键步骤。Radon 测度一词保留原文。」
